% Nagłówek LaTeX dla prac wzorcowych praca-magisterska.pl (XeLaTeX)
% W odróżnieniu od .wzory-src/header.tex: to jest produkt PŁATNY, więc bez
% stopki „Wzór do adaptacji" i bez promo smart-edu na każdej stronie.
\usepackage{fontspec}
\setmainfont{Times New Roman}
\usepackage{polyglossia}
\setdefaultlanguage{polish}

% Układ strony zgodny z typowym regulaminem prac dyplomowych:
% margines wewnętrzny szerszy pod oprawę.
\usepackage{geometry}
\geometry{a4paper, top=2.5cm, bottom=2.5cm, left=3.5cm, right=2.5cm}
% Interlinia idzie przez -V linestretch=1.5 w build-prace.py, nie przez
% \onehalfspacing tutaj: tcolorbox (ładowany niżej pod notę) resetował je
% po cichu i skład spadł z 72 na 56 stron przy niezmienionej treści.
\usepackage{indentfirst}
\setlength{\parindent}{1.25cm}

\usepackage[table]{xcolor}
\definecolor{brand}{HTML}{4F46E5}
\definecolor{annoGray}{HTML}{555555}
\definecolor{ruleGray}{HTML}{CCCCCC}

% Żywa pagina: tytuł skrócony po lewej, numer strony po prawej.
\usepackage{fancyhdr}
\pagestyle{fancy}
\fancyhf{}
\renewcommand{\headrulewidth}{0.4pt}
\renewcommand{\footrulewidth}{0pt}
\fancyhead[L]{\footnotesize\textcolor{annoGray}{\nouppercase{\leftmark}}}
\fancyfoot[C]{\footnotesize\thepage}
% Strony otwierające rozdział też mają numerację, ale bez żywej paginy.
\fancypagestyle{plain}{\fancyhf{}\fancyfoot[C]{\footnotesize\thepage}%
  \renewcommand{\headrulewidth}{0pt}}

% Tytuły rozdziałów i podrozdziałów
\usepackage{titlesec}
\titleformat{\section}{\Large\bfseries}{}{0pt}{}
\titlespacing*{\section}{0pt}{0pt}{18pt}
\titleformat{\subsection}{\large\bfseries}{\thesubsection}{0.6em}{}
\titlespacing*{\subsection}{0pt}{16pt}{8pt}
\titleformat{\subsubsection}{\normalsize\bfseries}{\thesubsubsection}{0.6em}{}

% Tabele — czytelne, bez pionowych linii (styl booktabs)
\usepackage{booktabs}
\usepackage{longtable}
\usepackage{array}
\renewcommand{\arraystretch}{1.25}

% Podpisy tabel i rysunków
\usepackage[font=small,labelfont=bf,skip=6pt]{caption}
% Praca numeruje tabele sama, w konwencji rozdziałowej („Tabela 1.1.”), i tak są
% wypisane w jej własnym spisie tabel. LaTeX dokładał do tego swoją numerację
% ciągłą, więc podpis brzmiał „Tabela 1: Tabela 1.1. Klasyfikacja…”.
\captionsetup[table]{labelformat=empty}
\captionsetup[longtable]{labelformat=empty}

% Bibliografia jako lista numerowana bez dodatkowych odstępów
\usepackage{enumitem}
\setlist[enumerate]{itemsep=3pt, parsep=0pt}

% Odsyłacze do bibliografii w kolorze marki, reszta czarna
\usepackage{hyperref}
\hypersetup{
  colorlinks=true,
  linkcolor=black,
  citecolor=brand,
  urlcolor=brand,
  linktoc=all,
  pdfborder={0 0 0}
}

% Streszczenie / abstrakt — węższy blok, mniejszy stopień pisma
\newenvironment{streszczenie}
  {\begin{quote}\small}
  {\end{quote}}

% Nota o charakterze danych w rozdziale badawczym.
% Rozdział empiryczny każdej pracy zawiera wyniki przykładowe, nie pochodzące
% z przeprowadzonego badania. Kupujący ma to wiedzieć z dokumentu, a nie
% dopiero z regulaminu — stąd ramka na starcie rozdziału i akapit na stronie
% redakcyjnej.
% Znaki, których Times New Roman nie ma. Bez tego xelatex wypisuje
% „Missing character" i wyrzuca je z PDF-u bez śladu — strzałka w schemacie
% przepływu znika, a zdanie przestaje się kleić.
\usepackage{newunicodechar}
\newunicodechar{→}{$\rightarrow$}
\newunicodechar{⟶}{$\longrightarrow$}
\newunicodechar{⟵}{$\longleftarrow$}
\newunicodechar{↔}{$\leftrightarrow$}
\newunicodechar{✔}{$\checkmark$}
\newunicodechar{✓}{$\checkmark$}
\newunicodechar{⚠}{\textbf{!}}

\usepackage{tcolorbox}
\definecolor{notaBg}{HTML}{FFF7ED}
\definecolor{notaFrame}{HTML}{EA580C}
\newtcolorbox{nota}{%
  colback=notaBg, colframe=notaFrame, boxrule=1pt, arc=2mm,
  left=10pt, right=10pt, top=9pt, bottom=9pt,
  before skip=12pt, after skip=14pt}
